% Pursuing Stacks (cotes 134-2 to 134-8): archive cote 134-7, site page 15 (PDF page 16), typescript page 457.
% typescript. Draft compared: ps6.tex:691-749.
% AI-transcribed by Claude Opus 5.5 (claude-opus-5-5, adaptive thinking, effort high), Message Batches API, batch ps3 (msgbatch_01XdFmuVVBAcgc7EFxdtCqsj), 2026-11-12; not yet checked by eye.
\scanpage{134-7/15}
\mspage{457}

It may be remarked that if M is any k-module,  than the property
that M be projective,  or flat,  can be described in terms of the isomor-%
phism class of the corresponding object X in U(k)\hand{$^\bullet$},  or equivalently, of
the functor W(M) on $\mathrm{Alg}_{/k}$ ,  with values in the category of pointed sets.
Indeed, the isomorphism class of the k-module $\Gamma_k(M)$  depends only on the
class of X,  and it is easily seen that M is projective, resp. flat, iff
$\Gamma_k(M)$ is.  (The ``only if'' is standard knowledge of commutative algebra,
the ``if'' comes from the fact that M is a direct factor of $\Gamma_k(M)$,  hence
projective resp. flat if the latter is.)  Using this,  one even checks
that it is enough to know the isomorphism class of X in U(k)  -  because
two objects of $U(k)^\bullet$  are isomorphic in $U(k)^\bullet$  iff they are in U(k).

* \quad *

*

The second afterthought is  that in the question on page 452,  we
should definitely assume $k=\underline{Z}$ .  Already when asserting hastily (p.453)
that an Eilenberg-Maclane homology group $H_i(\hand{$\pi$},n;\ \underline{Z})$  should automatically
inherit a structure of a k-module,  whenever \hand{$\pi$} had one,  I was feeling
unconfortable,  because after all the dependance of this group upon
variable \hand{$\pi$} is not at all additive,  thus the operation of k upon this
group stemming from it's operations on \hand{$\pi$} \xout{did} \unsure{were|was} not too likely to come
out additive !  To take one example,  take n=1 i.e. we just take ordina-%
ry group homology for \hand{$\pi$} ,  and assume that k is free of finite type over
$\underline{Z}$ ,  and \hand{$\pi$}\inserted{$=M$} free of finite type over k (for instance \hand{$M$}= k),  hence free
of finite type over $\underline{Z}$ too.  Then it is well known that
\[ H_*(M;\ \underline{Z})\hand{$\simeq \bigwedge$}_{\underline{Z}}M \qquad , \]
the exterior algebra of M  over $\underline{Z}$ ,  which surely is not endowed with a
structure of a k-module in any natural way !  Thus,  if there was any
such structure (natural or not) on the highest non zero term (correspon-%
ding to the \inserted[typed]{rank d}\xout{rak} of M over $\underline{Z}$ )\hand{,} it would follow that we get a ring homomor-%
phism  from  k to  $\underline{Z}\simeq \mathrm{End}\hand{$_{\underline{Z}}$}(H_d\simeq \underline{Z})$ ,  and \inserted[typed]{we}\xout{there} may choose k in such a
way that there is no such homomorphism,  for instance  $k = \underline{Z}[T]/(T^2+1)$ .
In this case,  there cannot be \uline{any} isomorphism  between the $H_d$'s of the
two members of \hand{(}u)\over{.}{!}  Another point,  which I hit upon first,  is that
when  $X_*$  is defined in terms of a semisimplicial object $M_*$ of $\mathrm{Ab}_k$ ,
then the functorial dependance of the first term in\hand{(w}) with respect to

\margin{15}
